\begin{longtable}{p{0.28\linewidth} p{0.62\linewidth}}

\textbf{Albedo}
  & The intrinsic material colour of a surface, independent of the illumination.
    A white wall has high albedo in all three colour channels; a black wall has
    low albedo.  In the IID model $\mathbf{I} = \mathbf{A} \cdot S$, the albedo
    is the term $\mathbf{A}$. \\[4pt]

\textbf{Chromaticity}
  & The colour of a pixel or region expressed independently of its overall
    brightness.  For an RGB pixel $(r, g, b)$, the chromaticity is typically
    the normalised vector $(r, g, b) / (r + g + b)$.  A constancy metric
    based on chromaticity measures whether the hue of the predicted albedo
    shifts when the illuminant changes. \\[4pt]

\textbf{Colour cast}
  & A global or regional shift in the hue of an image caused by a coloured
    light source.  A warm orange lamp produces an orange cast; a fluorescent
    tube produces a green cast.  When a cast leaks into the predicted albedo
    it means the model has not separated the illuminant colour from the
    material colour. \\[4pt]

\textbf{Illuminant}
  & The light source or combination of light sources illuminating a scene.
    An illuminant is characterised by its spectral power distribution, which
    determines both its brightness and its colour (e.g.\ warm, cool, neutral). \\[4pt]

\textbf{Intrinsic Image Decomposition (IID)}
  & The task of decomposing a photograph into an albedo layer (material colour)
    and a shading layer (lighting effect).  See Chapter~\ref{ch:method} for
    the formal definition. \\[4pt]

\textbf{Diffuse shading}
  & The per-pixel quantity that encodes how much light falls on each surface
    point, combining the effect of surface normal orientation, light distance
    and direction, and occlusion (shadows).  In this work the shading is a
    three-channel (colored) quantity $\mathbf{S}_d$, so it can represent a
    coloured illuminant; see Chapter~\ref{ch:method}. \\[4pt]

\textbf{Non-diffuse residual}
  & The remaining image content not explained by the albedo-times-diffuse-shading
    product: specular highlights, visible light sources, and clipped regions.
    Computed analytically as $\mathbf{R} = (\mathbf{I} - \mathbf{A} \odot \mathbf{S}_d)_{+}$,
    with no dedicated network head. \\[4pt]

\textbf{White-balance (WB)}
  & An image processing step that neutralises the colour of the illuminant,
    making images appear as if they were lit by a pure white light.
    White balance is applied by cameras automatically or in post-processing.
    In this work, WB is intentionally \emph{disabled} for the MID training
    pairs because it removes the illuminant colour information that the CARI
    training strategy relies on. \\[4pt]

\textbf{Scale-and-shift invariance (SSI)}
  & A property of a metric or loss that makes it insensitive to global
    multiplicative and additive offsets.  An SSI shading loss aligns the
    predicted shading to the ground truth using the best global scale
    and bias before computing the error, so it measures the quality of
    the relative shading pattern rather than its absolute level. \\[4pt]

\end{longtable}
